\section{Introduction}
\label{sec:intro}

\subsection{The paradox is an inference problem}

The universe is old, planets are common, and many planetary systems predate the
Solar System by billions of years \citep{lineweaver2001,borucki2016}. Yet no
extraterrestrial signal, probe, artefact or unambiguous technosignature has been
confirmed. The tension between these statements is conventionally called the
Fermi paradox, and is usually compressed into the question ``where is
everybody?'' \citep{hart1975,brin1983silence,webb2015}.

Stated that way the question is misleading, because it treats non-observation as
evidence about existence. It is not. Non-observation constrains the
\emph{product} of several quantities: how many civilizations arise, how many
survive to a technologically active phase, how strong a signature they produce,
and whether our searches have covered the region of parameter space in which
that signature would appear \citep{gray2015,tarter2001,sheikh2020}. A more
careful formulation asks which assumptions would make extraterrestrial
technological civilizations visible \emph{to us, now}.

This distinction is not merely rhetorical. It determines what a model of the
paradox should vary. Most quantitative treatments --- from the Drake equation
onward --- vary the number of civilizations while holding detectability fixed,
typically inside a single lifetime parameter $L$
\citep{prantzos2013,prantzos2020}. The alternative pursued here is to give
detectability its own dynamics.

\subsection{Existence, detectability, observation}

We distinguish three quantities that are frequently conflated:
\begin{equation}
  \Ntrue \;\geq\; \Ndet \;\geq\; \Nobs ,
  \label{eq:threeN}
\end{equation}
the number of civilizations that exist, the number producing signatures
detectable in principle, and the number actually detected by our searches. A
civilization may exist without being detectable; it may be detectable in
principle without ever being observed, if our search has not intersected its
signature. \Established

The framework developed here modifies the middle term. It asks how a
civilization's external observability evolves as its capability, regulatory
capacity, energy use, expansion behaviour and information efficiency change ---
and in particular what happens across a transition to recursive, self-improving
technological capability.

\subsection{Why recursive capability is the interesting transition}

A civilization becomes qualitatively different when its tools improve the process
of making better tools. Automated discovery, engineering and design close a loop
between capability and the rate of capability growth \citep{good1965,bostrom2014,russell2019}.
Quantum computation enters this picture in exactly this way: not as a separate
phenomenon, but as one possible contributor to the rate at which a civilization
can recursively improve its own technology, by accelerating the simulation,
search and optimisation problems that gate technological progress
\citep{shor1994,nielsen2010,preskill2018}. \Speculative

Such a loop plausibly changes more than growth rate. It may change information
efficiency, and hence signal leakage; it may change whether expansion is
attractive; it may shorten or lengthen the interval over which a civilization is
loud. None of this is asserted here as fact. The model asks what \emph{follows}
if such feedback occurs, and makes the dependence on that assumption explicit.

\subsection{What breaks if observability is modelled carelessly}

Giving observability its own differential equation is easy to do badly, and the
failure mode is instructive enough that we report it. If the terms representing
compression, efficiency and deliberate quietness are written as an absolute
subtraction --- a flux removed from $\Obs$ at a rate that does not depend on
$\Obs$ --- then nothing prevents observability from passing through zero and
running to large negative values. A trajectory can then be classified as
``quiet'' on the strength of a physically meaningless number.

That is not a coding slip; it is a modelling error with a clear physical
diagnosis. Compression and stealth reduce the signal a civilization
\emph{emits}. They are fractional reductions, and one cannot suppress emissions
that are not being made. Writing them as fractional rates acting on the emitted
signal, together with an explicit floor representing irreducible waste heat,
makes observability bounded below by construction. We prove this as
Proposition~\ref{prop:obs} and give the analogous statement for regulatory
capacity as Proposition~\ref{prop:reg}.

The correction has a substantive consequence, which is the main modelling result
of this paper. Once suppression is fractional, the quasi-steady observability is
a \emph{ratio} of production to suppression. Peak-then-decline behaviour then
follows automatically whenever suppression grows faster in capability than
production does --- it no longer has to be imposed by selecting a peaked
functional form. A low-observability account of the Great Silence that assumes a
peaked $\Obs(\Icap)$ assumes its conclusion; one that derives the peak from the
ratio of two independently motivated mechanisms does not.

\subsection{Contributions}

\begin{enumerate}\itemsep2pt
  \item An observability-aware decomposition
    $\Nobs = \Ntrue\,\Psurv\,h(\Obs)\,\Psearch$ that separates survival, signal
    production and search coverage, and locates precisely where a dynamical
    detectability model acts on the Drake equation (\S\ref{sec:decomp}).
  \item A three-variable dynamical model of capability, regulation and
    observability derived from mechanism balances, with two structural
    guarantees --- $\Obs \geq \Ofloor$ and $\Reg \in [0,1]$ --- that hold by
    construction rather than by clamping (\S\ref{sec:model}).
  \item Peak-then-decline observability as a \emph{derived} consequence of
    suppression outgrowing production, rather than an assumed functional form
    (\S\ref{sec:results}).
  \item A registry architecture in which every functional form is
    interchangeable, making the question ``do the conclusions survive changing
    the functions?'' a menu selection rather than a rewrite (\S\ref{sec:model}).
  \item A verification protocol unusual in this literature: closed-form
    solutions, an independent high-order integrator, property-based testing over
    randomised parameter sets, and two independently written implementations
    agreeing to solver tolerance (\S\ref{sec:verification}).
  \item An explicitly negative result about the Lambert-$W$ threshold that
    motivated this work: the balance that justifies $W$ is not the balance
    governing the integrated system (\S\ref{sec:lambert}).
\end{enumerate}

\subsection{What this paper does not claim}

It does not claim that extraterrestrial civilizations exist, that advanced
civilizations become invisible, that recursive artificial intelligence causes
collapse or transcendence, or that the Fermi paradox is solved. Every
substantive statement carries a strength label --- \Established, \Plausible,
\Hypothetical{} or \Speculative{} --- and these are collected in
Table~\ref{tab:claims}.
